\scanpage{610}\mspage{19}
% archive sheet number: 609

$\underline S$ n'est pas accessible dans $M$), on ne peut donc
conclure de cette inaccessibilité de $\underline S$, celle de
$\underline S'$. Notons le

\uline{Corollaire} Soit $\pi$ un cardinal infini. Pour que la
sous-catégorie fermée $\underline S'$ de
$\deleted{\unsure{\underline S}}\unsure{M|\underline M}$ (formée
des \deleted{$A\in\mathrm{Ob}\,M$} $k$-alg. comm. $A$ telles que
$\mathrm{Hom}_M(A,k)=\emptyset$) soit $\pi$-accessible, il f et s
que pour tte $k$-algèbre $A$ [de cardinal $>\pi$],
\inserted{telle que $\mathrm{Hom}_M(A,k)=\emptyset$,}
\unsure{existe} une sous-algèbre $A_i$
\uline{de cardinal} $\leq\pi$, telle que
$\mathrm{Hom}_M(A_i,k)=\emptyset$.

(On applique le lemme, en prenant pour $\alpha$ \unsure{le} premier
ordinal infini.)

\uline{Question} Cela est-il vérifié, en prenant \inserted{\unsure{simplement|simplemt}}
$\pi\geq\mathrm{Card}\,k$ ? La question de l'existence
d'un cardinal $\pi$, ayant la propriété indiquée plus haut me paraît
délicate, et j'ignore la réponse. Si on veut construire un
contre exemple pour tt $\pi$, les développements précédents suggèrent
de prendre pour $A$ un produit infini d'extensions non triviales de
$k$, ou un \deleted{\unsure{corps}}\inserted{\unsure{anneau}} $k_{\Phi}$ comme dans
la dém. de prop. 4 a).
